\section{Related work}
\label{sec:related-work}

\subsection{Formal semantics and verified subsets}

Although the syntax of mainstream languages is usually described precisely, few define their behaviour
formally, so well-defined subsets have been constructed to equip languages with formal, often mechanised,
semantics. Examples include Java~\cite{igarashi2001}, OCaml~\cite{seassau2025},
JavaScript~\cite{bodin2014,guha2010,park2015,politz2012}, PHP~\cite{filaretti2014},
Rust~\cite{jung2020,jung2018}, and lower-level languages such as WebAssembly~\cite{watt2021a,watt2021,watt2023}
and LLVM-IR~\cite{zakowski2021,zhao2012}. A notable success is Clight~\cite{blazy2009}, a substantial subset of
C that underpins a verified compiler~\cite{leroy2009}; likewise the CakeML project~\cite{kumar2014} produced a
verified implementation of a fragment of ML~\cite{milner1997}, building on earlier mechanisations~\cite{lee2007}.

\subsection{Formal semantics of Python}

Several efforts formalise Python itself: an embedding~\cite{guth2013} in the $\mathbb{K}$
framework~\cite{rosu2010}, a structural operational semantics~\cite{kohl2021}, a reduction of much of the
language to a core calculus~\cite{politz2013}, an executable semantics aimed at optimising
compilers~\cite{melancon2023}, and a skeletal semantics of a fragment~\cite{andrieux}. \PurePy differs in
restricting the language to a well-behaved subset rather than modelling Python in full; relating our semantics
to these is future work (\secref{agreement}).

\subsection{Language subsets for a purpose}

A separate line of work delineates subsets for a particular use. The \emph{Source} family~\cite{anderson2021}
comprises progressively larger subsets of JavaScript for teaching first-year students the foundations of
computing, following a JavaScript edition of a classic text~\cite{abelson2022}; the choice of JavaScript was
pragmatic, since industrial relevance weighs heavily in language choice for introductory
courses~\cite{simon2018}. For Python, ChocoPy~\cite{padhye2019} is a statically-typed subset used as the source
language in compiler-construction courses, and Restricted Python (RPython)~\cite{ancona2007} is a statically
typed subset that the PyPy project~\cite{rigo2006} compiles to common runtimes.

\subsection{Functional Python}

We are aware of one other functional subset of Python~\cite{sunderraman2025}, intended for teaching; \PurePy
encloses more of the language, since that work omits control structures such as conditionals. Related efforts
sit outside the subset discipline: Coconut\footnote{\url{https://coconut-lang.org}} is a functional language that compiles to Python, and
functional pipelines have been proposed for scientific computation in Python~\cite{zhang2024}. Although \PurePy does not aim to make Python
more functional, the community has moved in that direction, from lambda expressions in Python~1.0 to a recent
proposal for deeply immutable objects~\cite{johnson2025}. Typed treatments of effects in functional languages,
such as Prism,\footnote{\url{https://www.stephendiehl.com/posts/prism/}} are relevant to the structured effects \PurePy may add in a later release.

\subsection{Pattern matching}

Checking that the cases of a match are reachable and exhaustive goes back to the usefulness algorithm for
ML~\cite{maranget2007}, which asks of each case whether some value reaches it, and to the term descriptions a
match compiler carries as it decides a pattern one position at a time~\cite{sestoft1996}. \PurePy computes
both from a residual, the set of values no earlier case has matched, and seeds that residual from the type of
the scrutinee, so the types of the bindings a pattern makes come from the same computation. Accounting for
the type in this way follows the uncovered sets of \citet{karachalias2015}, which track what a GADT match
leaves and rule out cases the indices make impossible.
